\chapter{Lead--Lag et trading inter-plateformes}\label{ch:hf:lead-lag-and-cross-venue-trading}

Quand le future sur indice bouge, le fonds suit en quelques millisecondes et les actions qui le composent un peu plus tard. Qui mène dépend de l'heure, de la
plateforme et de l'ampleur du mouvement, et cela a changé à mesure que les firmes les plus rapides ont rattrapé leur retard. Entre 2005 et 2011, la durée médiane d'un arbitrage entre
le future S\&P 500 et son fonds est tombée de 97 millisecondes à 7. Dans les marchés simulés de ce chapitre, où un fonds suit son future avec une demi-seconde
de retard, l'avance est retrouvée à partir des deux mids à 25 millisecondes près. La part d'information du future est de 97--98~\%. Un preneur de liquidité qui négocie le fonds quand
le future a bougé et que le fonds n'a pas encore suivi gagne 0.49 tick par action s'il agit dans la milliseconde. L'edge a disparu avant la fin de la demi-seconde.

\section{Formation des prix entre marchés}

Deux prix d'une même chose (un future et le fonds qui réplique le même indice, une action négociée sur deux plateformes, une cryptomonnaie échangée sur deux bourses) partagent un
prix efficient commun et ne diffèrent que par un bruit transitoire. La formation des prix est le processus par lequel une nouvelle information entre dans ce prix commun~; le marché où elle
arrive en premier mène, et l'autre suit. Hasbrouck a établi, pour les indices actions américains, que l'essentiel de la formation des prix se fait sur les E-mini futures du S\&P
500 et du Nasdaq-100, que cette formation se partage entre le future standard et le fonds du S\&P 400, et que le fonds S\&P 500 mène les fonds sectoriels
bien plus qu'ils ne le mènent. Mizrach et Neely ont établi que les futures sur Treasuries portent l'essentiel de la formation des prix des obligations à cinq et dix ans à partir de 1999.

Aux fréquences les plus élevées, les deux prix ne bougent même pas ensemble. Budish, Cramton et Shim ont mesuré, pour le future S\&P 500 et son fonds en 2011,
une corrélation médiane des rendements de 0.10 à 10 millisecondes et de 0.008 à une milliseconde, contre près de 1 sur une minute~: l'effet Epps
(livre~4, chapitre~21) se transforme en une opportunité pour quiconque lit le meneur en premier.

\begin{definition}[Avance inter-plateformes]\label{def:hf:lead-lag-and-cross-venue-trading:lead}
Une \emph{avance inter-plateformes}\index{avance inter-plateformes} est le délai dont les prix d'un marché (une plateforme, ou un instrument partageant le même
prix efficient) devancent ceux d'un autre~: le décalage pour lequel la corrélation croisée de leurs variations de prix atteint son maximum, estimé sur leurs propres
ticks asynchrones.
\end{definition}

\section{Parts d'information}

Le livre~10, chapitre~8 définit les deux mesures de la formation des prix. La \emph{part d'information} (information share) d'un marché est la part de la variance des
innovations du prix efficient commun qui proviennent des innovations de ce marché~; la \emph{part de composante} (component share) est le poids de ce marché dans le
prix efficient commun. Les deux proviennent d'un modèle vectoriel à correction d'erreur (VECM) (livre~4, chapitre~20) portant sur les deux prix $p^1,p^2$ sur une grille régulière,
avec le vecteur de cointégration fixé à $(1,-1)$~:
\[
\Delta p_t=\alpha\,(p^1_{t-1}-p^2_{t-1})+\sum_{i=1}^{k}\Gamma_i\,\Delta p_{t-i}+e_t .
\]
Le marché qui se recale sur l'autre a un grand coefficient d'ajustement~; celui du meneur est proche de zéro.

\begin{proposition}[Les deux parts d'une paire]\label{prop:hf:lead-lag-and-cross-venue-trading:shares}
Avec $\alpha=(\alpha_1,\alpha_2)$ et $\alpha_\perp=(\alpha_2,-\alpha_1)$, la part de composante du marché 1 vaut $\gamma_1=\alpha_2/(\alpha_2-\alpha_1)$. Avec
$\Omega=\mathrm{Cov}(e_t)$ et $F$ son facteur de Cholesky pour un ordre donné des marchés, la part d'information du marché 1 est
$[(\gamma^\top F)_1]^2/\gamma^\top\Omega\gamma$, et les deux ordres en donnent les bornes inférieure et supérieure.
\end{proposition}

\admitted
La démonstration de la représentation à tendance commune est celle de Hasbrouck et de Gonzalo et Granger~; la forme ci-dessus utilise le fait que, pour une paire dont
le vecteur de cointégration $(1,-1)$, que les poids de la tendance commune sont proportionnels à $\alpha_\perp$.

La \Cref{fig:hf:lead-lag-and-cross-venue-trading:shares} montre les parts du meneur sur les marchés du chapitre. Pour une avance de 50 millisecondes, plus courte
que la grille de 0.1 seconde, les innovations des deux marchés sont simultanées sur la grille et la part d'information n'est pas identifiée~: ses bornes valent
0.52 et 0.73. À une demi-seconde, elles valent 0.97 et 0.98~; à une seconde, les deux valent 0.998. Les bornes se resserrent à mesure que la grille résout l'avance, ce qui explique
pourquoi les parts d'information calculées sur des données à la seconde disent peu de chose d'une course à la milliseconde.

\begin{omfigure}
\begin{tikzpicture}
\begin{axis}[width=11cm,height=5.2cm,xlabel={\small avance implantée de A sur B (secondes)},ylabel={\small part de A},xmode=log,log basis x=10,
  xtick={0.05,0.2,0.5,1},xticklabels={0.05,0.2,0.5,1},xmin=0.04,xmax=1.25,ymin=0.4,ymax=1.02,tick label style={font=\footnotesize},
  grid=major,grid style={gray!25},table/col sep=comma,legend style={font=\scriptsize,at={(0.5,-0.3)},anchor=north},legend columns=3]
\addplot[omDef,thick,mark=*,mark size=1.3pt] table[x=lag,y=is_high]{figdata/hft/08-lead-lag-and-cross-venue-trading/shares.csv};
\addplot[omDef,thick,dashed,mark=o,mark size=1.3pt] table[x=lag,y=is_low]{figdata/hft/08-lead-lag-and-cross-venue-trading/shares.csv};
\addplot[omThm,thick,mark=square*,mark size=1.3pt] table[x=lag,y=component]{figdata/hft/08-lead-lag-and-cross-venue-trading/shares.csv};
\legend{part d'information (borne haute),part d'information (borne basse),part de composante}
\end{axis}
\end{tikzpicture}

\omcaption{Parts de formation des prix du marché meneur A, issues d'un modèle vectoriel à correction d'erreur sur une grille de 0.1 seconde avec dix retards, selon l'avance
implantée dans la simulation (échelle logarithmique)~; moyenne de deux sessions de vingt minutes. Données~: \texttt{hf\_leadlag.table}.}\label{fig:hf:lead-lag-and-cross-venue-trading:shares}
\end{omfigure}

\section{Estimer le décalage sur des ticks asynchrones}

Les prix de deux marchés changent à des moments différents~; échantillonner les deux sur une grille ajoute le biais d'Epps (livre~4, chapitre~21). L'estimateur d'avance--retard
du livre~7, chapitre~10 travaille sur les ticks eux-mêmes~: la covariance de Hayashi--Yoshida des accroissements des deux séries, l'horloge de l'une
translatée d'un décalage candidat, normalisée en corrélation~; le décalage qui donne la plus forte corrélation est l'avance
(\cref{fig:hf:lead-lag-and-cross-venue-trading:ccf}).

\begin{omfigure}
\begin{tikzpicture}
\begin{axis}[width=11cm,height=5.2cm,xlabel={\small décalage candidat de B derrière A (secondes)},ylabel={\small corrélation croisée},
  tick label style={font=\footnotesize},xmin=-1.5,xmax=1.5,ymax=0.4,grid=major,grid style={gray!25},table/col sep=comma]
\addplot[omDef,thick] table[x=lag,y=ccf]{figdata/hft/08-lead-lag-and-cross-venue-trading/ccf.csv};
\draw[omThm,dashed] (axis cs:0.5,0) -- (axis cs:0.5,0.4);
\node[font=\scriptsize,omThm,anchor=north east,fill=white,inner sep=1pt] at (axis cs:0.4,0.39) {avance implantée 0.5~s};
\end{axis}
\end{tikzpicture}

\omcaption{Corrélation croisée de Hayashi--Yoshida des variations des deux mids, sur leurs propres ticks, l'horloge de B étant reculée de chaque décalage candidat, en
pas de 25 millisecondes~; une session de vingt minutes avec une avance implantée d'une demi-seconde. Données~: \texttt{hf\_leadlag.estimate}.}\label{fig:hf:lead-lag-and-cross-venue-trading:ccf}
\end{omfigure}

Sur deux sessions pour chaque avance implantée, les estimations valent 0.05 et 0.25 secondes pour 0.05, 0.225 et 0.25 pour 0.2, 0.475 et 0.575 pour 0.5, et 1.075
et 1.1 pour une seconde. Un marché peut être trop lent pour cela~: dans la configuration par défaut du livre~7, chaque carnet simulé suit son
prix efficient sur plusieurs secondes, le maximum de la corrélation croisée est plat, et le même estimateur a renvoyé des avances de $-1.65$ et $+1.0$ secondes pour une avance implantée de 50
millisecondes. Les marchés du chapitre rendent les traders informés quatre fois plus actifs et les cotations périmées s'annulent plus vite, si bien que chaque mid suit son
prix efficient à une fraction de seconde près~; ce n'est qu'alors qu'une avance inférieure à la seconde devient identifiable.

\section{Trader l'avance~: coter et prendre}

Un teneur de marché sur le suiveur utilise l'avance pour protéger ses cotations (le skew piloté par le future, vu au chapitre~7). Un preneur de liquidité l'utilise pour franchir la fourchette du suiveur
avant que le suiveur ne bouge. La règle du preneur est ici la suivante~: lorsque le mid du meneur a bougé d'au moins un tick sur la dernière seconde de plus que celui du suiveur,
franchir la fourchette du suiveur pour un lot, et valoriser la transaction au mid du suiveur cinq secondes plus tard, une fois déduits une demi-fourchette et des frais de 0.1 tick.
Agir $d$ secondes après le mouvement du meneur (voir la \cref{fig:hf:lead-lag-and-cross-venue-trading:edge}~; avance implantée d'une demi-seconde, 135 transactions sur deux
sessions)~:

\begin{omfigure}
\begin{tikzpicture}
\begin{axis}[width=11cm,height=5.2cm,xlabel={\small latence du preneur après le mouvement du meneur (millisecondes)},ylabel={\small edge (ticks par action)},
  xmode=log,log basis x=10,xmin=0.7,xmax=1400,xtick={1,10,100,1000},xticklabels={1,10,100,1000},ymin=-0.8,ymax=0.8,
  tick label style={font=\footnotesize},grid=major,grid style={gray!25},table/col sep=comma]
\addplot[omThm,thick,mark=*,mark size=1.3pt,error bars/.cd,y dir=both,y explicit] table[x=latency_ms,y=edge,y error=se]{figdata/hft/08-lead-lag-and-cross-venue-trading/edge.csv};
\draw[gray] (axis cs:0.7,0) -- (axis cs:1400,0);
\draw[omDef,dashed] (axis cs:500,-0.8) -- (axis cs:500,0.8);
\node[font=\scriptsize,omDef,anchor=north east,fill=white,inner sep=1pt] at (axis cs:480,0.75) {avance implantée};
\end{axis}
\end{tikzpicture}

\omcaption{Edge par action de la prise d'un lot du suiveur après que le meneur a bougé et que le suiveur ne l'a pas fait, une fois déduites la moitié de la fourchette et des
frais de 0.1-tick, valorisée cinq secondes plus tard, en fonction de la latence du preneur (échelle logarithmique), $\pm$ une erreur type~; avance implantée d'une demi-seconde, 135 transactions en deux
sessions de vingt minutes. Données~: \texttt{hf\_leadlag.edge\_curve}.}\label{fig:hf:lead-lag-and-cross-venue-trading:edge}
\end{omfigure}

L'edge est de 0.485 tick par action à une milliseconde, de 0.43 à 100, de 0.24 à 300, et de $-0.22$ à 500 millisecondes, l'avance implantée~: le suiveur a
bougé, l'écart a disparu, et le preneur paie la moitié de la fourchette et les frais pour rien ($-0.54$ à une seconde). La majeure partie de l'edge survit aux premières 100
millisecondes seulement parce que l'avance de ce marché est longue. Les chiffres de Budish, Cramton et Shim placent l'avance réelle entre le future et le fonds à quelques millisecondes en 2011,
Ainsi la même courbe, tracée pour une paire réelle, se comprime dans les premières millisecondes, et la vitesse décide qui la gagne (chapitre~9).

\section{Quand le meneur change}

Une avance est une propriété d'une paire à un instant donné. La plateforme qui mène peut changer selon l'heure (un marché de contrats à terme qui ouvre avant les actions mène au
moment de l'ouverture), selon le mouvement (les grands mouvements commencent sur le marché le plus profond), selon les frais ou les pannes d'une plateforme, et selon la concurrence~: quand assez de traders négocient
l'avance, les teneurs de marché du suiveur retirent plus vite leurs cotations périmées (chapitre~6), et l'avance se raccourcit. Les conséquences pratiques sont d'estimer
l'avance en continu sur les données récentes, selon l'heure de la journée, et de surveiller la part d'information pour un changement de meneur, comme pour toute
décroissance d'un prédicteur (le livre~7, chapitre~13).

\section{Fiches de stratégie}

\begin{strategyfile}{Future sur indice vers fonds}\label{strat:hf:lead-lag-and-cross-venue-trading:future}
\sfield{Qui paie, et pourquoi} Les fournisseurs de liquidité du fonds dont les cotations sont périmées pendant les millisecondes qui suivent le mouvement du future.
\sfield{Instruments et plateformes} Un future sur indice et le fonds indiciel coté (ETF) sur le même indice (et ses options).
\sfield{Signal} Le mouvement du future, exprimé en unités du fonds par la base de fair value (chapitre~2), moins le mouvement du fonds sur une courte fenêtre.
\sfield{Dimensionnement et exécution} Prendre une tranche du fonds lorsque l'écart dépasse la moitié de la fourchette plus les frais~; ou, en tant que teneur de marché du fonds, retirer le
côté périmé.
\sfield{Coûts} La moitié de la fourchette, les frais du preneur, la couverture dans le future si la position est conservée.
\sfield{Comment elle s'éteint} La course à la vitesse~: l'opportunité médiane a duré 97 millisecondes en 2005 et 7 en 2011, tandis que son profit médian par opportunité
est resté proche de 0.08 points d'indice (Budish, Cramton et Shim)~; l'edge revient au plus rapide.
\sfield{Horizon, capacité, infrastructure} Des millisecondes~; une tranche par opportunité~; la liaison la plus rapide entre les
moteurs d'appariement du future et du fonds (le livre~14).
\sfield{Backtest honnête} Les horodatages des deux plateformes sur une même horloge, le délai de propagation réel, et l'exécution au prix disponible au moment où l'ordre
arrive.
\sfield{Sources} Hasbrouck (2003)~; Budish, Cramton et Shim (2015)~; ce chapitre~: 0.485 tick par action à une milliseconde, disparu avant l'avance implantée.
\end{strategyfile}

\begin{strategyfile}{Du fonds vers ses composantes}\label{strat:hf:lead-lag-and-cross-venue-trading:etf}
\sfield{Qui paie, et pourquoi} Les teneurs de marché sur actions individuelles dont les cotations retardent un mouvement généralisé du marché qui apparaît d'abord dans le fonds ou le future.
\sfield{Instruments et plateformes} Un fonds large ou sectoriel et ses plus grandes composantes.
\sfield{Signal} Le mouvement du fonds multiplié par le bêta de chaque action, moins le mouvement propre de l'action.
\sfield{Dimensionnement et exécution} Décaler ou retirer le côté périmé de nombreuses actions à la fois~; ne prendre qu'aux plus grands écarts.
\sfield{Coûts} De nombreuses petites positions et leurs frais~; une couverture dans le fonds.
\sfield{Comment elle s'éteint} Nouvelles propres à une action~: une action qui n'a pas suivi peut avoir sa raison propre. Les fonds sectoriels suivent davantage le fonds large qu'ils ne le mènent
(Hasbrouck), si bien que le sens compte.
\sfield{Horizon, capacité, infrastructure} De la milliseconde à la seconde~; un portefeuille de plusieurs centaines de titres.
\sfield{Backtest honnête} Bêtas à la date~; ne retirer les actions ayant eu des nouvelles dans la fenêtre que si le système en production le pouvait.
\sfield{Sources} Hasbrouck (2003).
\end{strategyfile}

\begin{strategyfile}{Treasury au comptant vers futures sur Treasury}\label{strat:hf:lead-lag-and-cross-venue-trading:treasury}
\sfield{Qui paie, et pourquoi} Les fournisseurs de liquidité sur le plus lent des deux, le bon au comptant ou son future, selon le côté qui mène à ce moment-là.
\sfield{Instruments et plateformes} Les bons on-the-run sur les plateformes interdealer (le livre~2, chapitre~4) et les futures Treasury.
\sfield{Signal} Le mouvement de rendement du meneur, converti par le facteur de concordance et la base (le livre~2, chapitre~6) en prix du suiveur.
\sfield{Dimensionnement et exécution} Prendre chez le suiveur~; ou décaler les cotations à cet endroit.
\sfield{Coûts} Frais de plateforme, mouvements propres de la base, financement si la position est conservée.
\sfield{Comment elle s'éteint} Le meneur change~: Mizrach et Neely ont constaté que les futures portent l'essentiel de la formation des prix pour les bons à cinq et dix ans à partir de 1999,
la part du marché au comptant variant selon les fourchettes, l'activité et la volatilité.
\sfield{Horizon, capacité, infrastructure} Quelques millisecondes autour des publications de données, quelques secondes sinon~; accès aux moteurs d'appariement des deux marchés.
\sfield{Backtest honnête} Estimer le meneur sur les seules données passées, période par période.
\sfield{Sources} Mizrach et Neely (2008).
\end{strategyfile}

\begin{strategyfile}{De la plateforme crypto menante vers la plateforme en retard}\label{strat:hf:lead-lag-and-cross-venue-trading:crypto}
\sfield{Qui paie, et pourquoi} Les teneurs de marché des bourses plus lentes ou plus petites, dont les carnets d'ordres suivent tardivement les plateformes qui forment les prix.
\sfield{Instruments et plateformes} Un même jeton sur plusieurs plateformes centralisées (le livre~3, chapitre~16), au comptant et perpétuels.
\sfield{Signal} Le mouvement de la plateforme menante moins celui du suiveur, avance estimée par paire de plateformes.
\sfield{Dimensionnement et exécution} Prendre chez le suiveur avec un inventaire préalablement positionné là-bas (le livre~3, chapitre~16)~; ou coter sur place avec un skew.
\sfield{Coûts} Frais du preneur en points de base, transferts pour rééquilibrer l'inventaire, limites de débit.
\sfield{Comment elle s'éteint} Quand les capitaux circulent librement, les écarts se résorbent~; Makarov et Schoar ont constaté des écarts bien plus grands entre pays qu'à l'intérieur de chacun, avec
des contrôles des capitaux qui les maintiennent ouverts.
\sfield{Horizon, capacité, infrastructure} De la milliseconde à la seconde~; serveurs proches de chaque plateforme (le livre~14).
\sfield{Backtest honnête} Les horodatages de chaque plateforme et la latence de leurs données jusqu'au trader, et non une bande consolidée.
\sfield{Sources} Makarov et Schoar (2020)~; chapitre~24.
\end{strategyfile}

\begin{strategyfile}{De la plateforme FX principale vers les plateformes secondaires}\label{strat:hf:lead-lag-and-cross-venue-trading:fx}
\sfield{Qui paie, et pourquoi} Les fournisseurs de liquidité sur les plateformes secondaires dont les flux retardent sur la plateforme principale (le livre~2, chapitre~14).
\sfield{Instruments et plateformes} Les principales paires de devises sur la plateforme principale, les autres plateformes et les flux vers les agrégateurs.
\sfield{Signal} Le mouvement de la plateforme principale moins celui de la plateforme, et la cohérence triangulaire entre paires.
\sfield{Dimensionnement et exécution} Prendre là où la plateforme le permet~; le last look (chapitre~9) permet au fournisseur de rejeter.
\sfield{Coûts} Frais de plateforme et rejets.
\sfield{Comment elle s'éteint} Le trading algorithmique lui-même~: Chaboud et ses coauteurs ont constaté qu'il a réduit les opportunités d'arbitrage triangulaire, principalement parce que des ordinateurs
jouent le rôle de preneurs de liquidité.
\sfield{Horizon, capacité, infrastructure} Des millisecondes~; les sites d'appariement de la plateforme principale (le livre~14, chapitre~24).
\sfield{Backtest honnête} Inclure les rejets sous last look et les temps de rétention.
\sfield{Sources} Chaboud, Chiquoine, Hjalmarsson et Vega (2014).
\end{strategyfile}

\section{Tutoriel~: qui bouge en premier}\label{tut:hf:lead-lag-and-cross-venue-trading:tut}

\begin{tutorial}
\textbf{Objectif.} Retrouver une avance implantée à partir de deux séries de prix asynchrones, mesurer les parts de formation des prix, et valoriser l'avance en fonction de la
latence. \textbf{État final~:} \cref{fig:hf:lead-lag-and-cross-venue-trading:shares,fig:hf:lead-lag-and-cross-venue-trading:ccf,fig:hf:lead-lag-and-cross-venue-trading:edge}.
\begin{tutsteps}
\item \textbf{Marchés.} \texttt{hf\_leadlag.pair(lag, seed)} simule le meneur et le suiveur, rendus réactifs.
\item \textbf{Avance.} \texttt{estimate} applique l'estimateur de Hayashi--Yoshida du livre~7 aux points de changement des deux prix milieux sur une grille de retards candidats.
\item \textbf{Parts.} \texttt{firm.xvenue.shares} ajuste le modèle vectoriel à correction d'erreur sur une grille au pas de 0.1~seconde et renvoie la part de composante et les
bornes de la part d'information (le \cref{lst:hf:lead-lag-and-cross-venue-trading:shares}).
\omcode{code/firm/xvenue/firm_xvenue.py}{46}{71}{Le modèle vectoriel à correction d'erreur et les deux parts de formation des prix.}\label{lst:hf:lead-lag-and-cross-venue-trading:shares}
\item \textbf{Edge.} \texttt{lead\_edge} rejoue les deux carnets d'ordres sans impact et prend position sur l'écart après une latence~; \texttt{edge\_curve()} balaie ces valeurs.
\omcode{code/firm/xvenue/firm_xvenue.py}{81}{97}{Prendre position sur l'écart après une latence, au prix du suiveur à ce moment-là~; valoriser cinq secondes plus tard.}\label{lst:hf:lead-lag-and-cross-venue-trading:edge}
\end{tutsteps}
\textbf{Pour aller plus loin.} Rendre l'avance aléatoire d'un mouvement à l'autre et observer l'élargissement du pic de la corrélation croisée~; donner au suiveur sa propre
information (une part de son prix efficient bouge en premier) et observer les parts d'information se rejoindre.
\end{tutorial}

\section{Construction~: la boîte à outils inter-plateformes}\label{bld:hf:lead-lag-and-cross-venue-trading:xvenue}

\begin{build}
\bfield{Rôle} Qui mène, de combien, et ce que vaut l'avance à une latence donnée.
\bfield{Interface} \texttt{sample(t, x, grid)}~; \texttt{lead(t1, x1, t2, x2, lags)}~; \texttt{shares(p1, p2, lags)} avec la part de composante et les
bornes de la part d'information~; \texttt{lead\_edge(lead\_t, lead\_mid, fol\_t, fol\_bid, fol\_ask, latency, move, window, H, fee)}.
\bfield{Règles} Les avances sont estimées sur les propres ticks des séries~; les parts sur une grille dont le pas est indiqué~; les edges après la moitié de la fourchette et les frais, un lot,
sans impact.
\bfield{Tests d'acceptation} \texttt{code/firm/xvenue/tests/}~: sur une marche aléatoire observée par un marché trois pas avant l'autre, les parts du meneur
dépassent 0.8 et s'inversent lorsque les marchés sont inversés~; une avance implantée sur une grille est retrouvée exactement~; l'edge d'une paire construite à la main à deux latences.
\bfield{Pour aller plus loin} Parts par heure de la journée~; une avance qui dépend de la taille du mouvement~; plusieurs suiveurs à la fois.
\end{build}

\begin{omsources}
\item J.~Hasbrouck, ``One security, many markets: determining the contributions to price discovery'', \emph{Journal of Finance} 50(4), 1995.
\item J.~Gonzalo et C.~Granger, ``Estimation of common long-memory components in cointegrated systems'', \emph{Journal of Business and Economic
Statistics} 13(1), 1995.
\item J.~Hasbrouck, ``Intraday price formation in U.S. equity index markets'', \emph{Journal of Finance} 58(6), 2003.
\item B.~Mizrach et C.~J.~Neely, ``Information shares in the US Treasury market'', \emph{Journal of Banking and Finance} 32(7), 2008.
\item I.~Makarov et A.~Schoar, ``Trading and arbitrage in cryptocurrency markets'', \emph{Journal of Financial Economics} 135(2), 2020.
\item A.~Chaboud, B.~Chiquoine, E.~Hjalmarsson et C.~Vega, ``Rise of the machines: algorithmic trading in the foreign exchange market'',
\emph{Journal of Finance} 69(5), 2014.
\item E.~Budish, P.~Cramton et J.~Shim, \emph{Quarterly Journal of Economics} 130(4), 2015.
\end{omsources}

\section{Exercices}

\begin{exercise}[$\star$]\label{exo:hf:lead-lag-and-cross-venue-trading:1}
Un ajustement à correction d'erreur donne $\alpha_1=-0.02$ et $\alpha_2=0.18$. Quel marché mène, et quelle est sa part de composante~?
\end{exercise}

\begin{exercise}[$\star$]\label{exo:hf:lead-lag-and-cross-venue-trading:2}
Un preneur de liquidité réalise 135 transactions en deux sessions de vingt minutes, à 0.363 tick par action, avec un lot de 100 actions et un tick d'un centime. Combien gagne-t-il
par session~?
\end{exercise}

\begin{exercise}[$\star$]\label{exo:hf:lead-lag-and-cross-venue-trading:3}
De quel facteur la durée de vie médiane de l'arbitrage future--fonds a-t-elle diminué entre 2005 et 2011~?
\end{exercise}

\begin{exercise}[$\star\star$]\label{exo:hf:lead-lag-and-cross-venue-trading:4}
Pourquoi les bornes de la part d'information sont-elles larges pour une avance de 50 millisecondes et étroites pour une seconde, sur une grille de 0.1 seconde~?
\end{exercise}

\begin{exercise}[$\star\star$]\label{exo:hf:lead-lag-and-cross-venue-trading:5}
Pourquoi l'edge devient-il négatif, plutôt que nul, dès que la latence dépasse l'avance~?
\end{exercise}

\begin{exercise}[$\star\star$]\label{exo:hf:lead-lag-and-cross-venue-trading:6}
D'après la \cref{fig:hf:lead-lag-and-cross-venue-trading:edge}, entre quelles latences l'edge s'annule-t-il~?
\end{exercise}

\begin{exercise}[$\star\star\star$]\label{exo:hf:lead-lag-and-cross-venue-trading:7}
\emph{Programmation.} Avec une avance implantée de 0.2 seconde, calculer l'edge à des latences d'une milliseconde et de 100 millisecondes (\texttt{edge(0.2, 0.001)},
\texttt{edge(0.2, 0.1)}).
\end{exercise}

\begin{exercise}[$\star\star\star$]\label{exo:hf:lead-lag-and-cross-venue-trading:8}
\emph{Trouver la faille.} «~Sur des données à une seconde, la part d'information du future est de 99~\%, donc le fonds ne contribue en rien à la formation des prix et
peut ignorer son carnet d'ordres.~»
\end{exercise}

\section{Problème~: qui bouge en premier}

\begin{problem}[{Problème du week-end --- qui bouge en premier}]\label{pb:hf:lead-lag-and-cross-venue-trading:1}
Un desk échange un fonds contre son future sur indice et doit décider à quelle vitesse il doit être.

\medskip\noindent\textbf{Partie I --- Les observations.}
\begin{enumerate}
\item Qu'a établi Hasbrouck sur la formation des prix dans les marchés d'indices actions américains~?
\item Qu'ont établi Mizrach et Neely sur les bons du Trésor~?
\item Que devient la corrélation future--fonds à 10 millisecondes et à 1 milliseconde~?
\item Définir une \omterm{def:hf:lead-lag-and-cross-venue-trading:lead}{avance inter-plateformes}.
\end{enumerate}

\medskip\noindent\textbf{Partie II --- Mesurer.}
\begin{enumerate}[resume]
\item Écrire le modèle à correction d'erreur et les deux parts de la \cref{prop:hf:lead-lag-and-cross-venue-trading:shares}.
\item Donner les parts du meneur pour chaque avance implantée et expliquer leur structure.
\item Pourquoi l'avance doit-elle être estimée sur les propres ticks des séries~?
\item Donner les estimations d'avance pour chaque avance implantée.
\item Pourquoi les carnets simulés par défaut masquaient-ils une avance de 50 millisecondes~?
\end{enumerate}

\medskip\noindent\textbf{Partie III --- Trading.}
\begin{enumerate}[resume]
\item Décrire la règle du preneur de liquidité et la manière dont son edge est mesuré.
\item Donner l'edge à 1, 100, 300 et 500 millisecondes, et à une seconde.
\item Pourquoi l'edge reste-t-il important à 100 millisecondes ici, et pourquoi ne le serait-il pas pour une vraie paire future--fonds~?
\item Comment un teneur de marché du suiveur utilise-t-il le même signal~?
\end{enumerate}

\medskip\noindent\textbf{Partie IV --- Le verdict.}
\begin{enumerate}[resume]
\item Énoncer le \emph{résultat nommé}~: l'avance implantée et estimée, la part d'information du meneur et la latence à laquelle l'edge disparaît.
\item Pourquoi le meneur change-t-il au fil du temps~?
\item Quelle fiche de stratégie dépend le plus du contrôle des capitaux, et pourquoi~?
\item Comment surveilleriez-vous un changement de meneur en production~?
\item Pourquoi l'edge ne dépend-il pas de la durée de l'avance tant que le preneur de liquidité est plus rapide qu'elle~?
\item Quel est le coût d'un avantage de latence qui n'est pas utilisé~?
\item En une phrase~: qu'est-ce qui détermine qui obtient une avance~?
\end{enumerate}
\end{problem}

\section{Questions d'entretien}

\begin{interviewq}[$\star$ \iqroles{trader}]\label{iq:hf:lead-lag-and-cross-venue-trading:1}
Le future monte de deux ticks et le fonds n'a pas bougé. Donnez trois raisons de ne pas acheter le fonds.
\end{interviewq}

\begin{interviewq}[$\star\star$ \iqroles{chercheur}]\label{iq:hf:lead-lag-and-cross-venue-trading:2}
Pourquoi l'échantillonnage de deux séries asynchrones sur une grille biaise-t-il leur corrélation à la baisse, et comment l'estimateur de Hayashi--Yoshida l'évite-t-il~?
\end{interviewq}

\begin{interviewq}[$\star\star$ \iqroles{chercheur}]\label{iq:hf:lead-lag-and-cross-venue-trading:3}
Expliquez pourquoi la part d'information admet une borne supérieure et une borne inférieure, et ce qui les rapproche.
\end{interviewq}

\begin{interviewq}[$\star\star$ \iqroles{développeur}]\label{iq:hf:lead-lag-and-cross-venue-trading:4}
Les flux de deux plateformes arrivent sur des serveurs situés dans deux centres de données. Comment les horodateriez-vous pour qu'une avance de quelques millisecondes puisse être mesurée~?
\end{interviewq}

\begin{interviewq}[$\star\star$ \iqroles{trader}]\label{iq:hf:lead-lag-and-cross-venue-trading:5}
Votre edge issu du trading d'une avance est divisé par deux en un mois, alors que sa fréquence est inchangée. Qu'est-ce qui s'est passé, et que vérifiez-vous~?
\end{interviewq}

\begin{interviewq}[$\star\star\star$ \iqroles{chercheur}]\label{iq:hf:lead-lag-and-cross-venue-trading:6}
Dans une paire cointégrée de vecteur $(1,-1)$, montrer que si le marché 2 corrige seul ($\alpha_1=0$), la part de composante du marché 1 vaut un, et discuter de sa
part d'information.
\end{interviewq}
